\documentclass[11pt,a4paper]{article}
\usepackage[T1]{fontenc}
\usepackage[utf8]{inputenc}
\usepackage{lmodern}
\usepackage[margin=25mm]{geometry}
\usepackage{microtype}
\usepackage{amsmath,amssymb,amsthm,mathtools,mathrsfs}
\usepackage{booktabs,array,longtable}
\usepackage{enumitem}
\usepackage{adjustbox}
\usepackage{url}
\usepackage[round,authoryear]{natbib}
\usepackage{xcolor}
\usepackage[colorlinks=true,linkcolor=blue!50!black,citecolor=blue!50!black,urlcolor=blue!60!black]{hyperref}
\usepackage[capitalise,noabbrev]{cleveref}
% Shared mathematical notation. Arc i -> j means: i requires j.
% Siamo nel mondo dei grafi e delle closure: V sono i vertici, non "item".
\newcommand{\Vset}{\mathcal V}
\newcommand{\Aset}{\mathcal A}
\newcommand{\Gset}{\mathcal G}
\newcommand{\Cset}{\mathscr C}
\newcommand{\Rset}{\mathbb R}
\newcommand{\Sseq}{\mathcal S}
\newcommand{\Fset}{\mathcal F}
\newcommand{\Qset}{\mathcal Q}
\newcommand{\Bset}{\mathcal B}
\newcommand{\Nset}{\mathcal N}
\newcommand{\Tset}{\mathcal T}
\newcommand{\Kset}{\mathcal K}   % blocchi canonici K_1, ..., K_k
% Matrici e vettori in grassetto; le lettere greche via \boldsymbol.
\newcommand{\mat}[1]{\mathbf{#1}}
\newcommand{\vect}[1]{\mathbf{#1}}
\newcommand{\vecg}[1]{\boldsymbol{#1}}
\newcommand{\ind}{\chi}
\newcommand{\divg}{\operatorname{div}}
\newcommand{\supp}{\operatorname{supp}}
% Nome del metodo (decisione D18): minuscolo, sul modello di *network simplex*.
% \FS in mezzo alla frase, \FSC a inizio frase, \CS per la forma breve.
\newcommand{\FS}{ratio-closure simplex}
\newcommand{\FSC}{Ratio-closure simplex}
\newcommand{\CS}{closure simplex}
\newcommand{\LP}{\textsc{LP}}
\newcommand{\rhoT}{\rho}

\newcommand{\file}[1]{\path{#1}}

\title{The ratio-closure simplex:\\ literature review and novelty positioning}
\author{Working document --- not a paper section}
\date{\today}

\begin{document}
\maketitle

\begin{abstract}
This document surveys the work that the \FS{} algorithm has to be positioned
against, and states what may and may not be claimed as new. It is deliberately
written to \emph{reduce} the claim: every section that identifies a classical or
recent antecedent is a section that removes novelty capital, and the positioning
in \cref{sec:positioning} is what survives. Claims that earlier drafts made and
that this review retracts are listed explicitly in \cref{sec:retracted}, because
a retraction that is not written down tends to reappear.

Every reference was checked against its source; the ones that could not be
obtained are listed in \file{literature/TO\_OBTAIN.md} rather than cited from
memory.
\end{abstract}

\tableofcontents

\section{Scope and method}\label{sec:scope}

\subsection{What is being positioned}

The object is the \FS{} algorithm of this project: an ordinary primal simplex
method applied to the normalised linear program of the maximum-ratio closure
oracle,
\[
\max\Bigl\{\textstyle\sum_i p_i y_i \;:\; \sum_i w_i y_i = 1,\;
y_i - y_j \le 0\ \ (i,j)\in\Aset,\; y \ge 0 \Bigr\},
\]
in which a basis is represented by a forest of tight precedence arcs together
with one distinguished component, and the outer procedure calls that oracle on
successive residual graphs to build the canonical sequence of blocks that
yields the LP value of the precedence-constrained knapsack at every capacity.

Four things therefore need positioning, and they have different literatures:
the closure model and its LP; the ratio objective; the residual decomposition
and its parametric reading; and the simplex-on-a-forest machinery itself. The
last is where the strongest antecedents sit, and \cref{sec:forests} is the
section that matters most.

\subsection{How the sources were handled}

Every reference in this document was checked against its source: journal or
publisher page for the classical items, arXiv abstract page for the preprints,
the original forum post for the one item that is a forum post. Metadata were not
reconstructed from memory. Items that could not be obtained in full text are
listed in \file{literature/TO\_OBTAIN.md} with what is needed and why, and are
cited here only for facts that the abstract or the record itself supports.

The full texts that were obtained sit in \file{literature/papers/}. Most come
from the project's own historical library in \file{PAPER\_LONG/BIBLIO}; the
recent ones were downloaded from open sources. Each file is named by its BibTeX
key --- \file{Yang2026.pdf} for \citet{Yang2026} --- so that a citation in this
document and the corresponding full text are one lookup apart, and so that the
two lists can be reconciled mechanically: \file{literature/check.sh} reports any
PDF without a bibliography entry and any entry without a PDF.
\file{literature/papers/README.md} is the index, with the provenance of each
file.

\subsection{A note on the direction of this review}

A literature review written by the authors of an algorithm has a systematic
bias, and the way to counter it is to make the review adversarial on purpose.
The organising question of every section below is therefore not \emph{what have
others done that is different}, but \emph{what have others already done that we
were about to claim}. \Cref{sec:retracted} lists what that question cost.

\section{Maximum closure and its linear program}\label{sec:closure}

A \emph{closure} of a directed graph is a set of vertices containing all successors
of its members; with the arc convention of this project, \(i \to j\) meaning
``\(i\) requires \(j\)'', a closure is a set closed under prerequisites. The
maximum-weight closure problem is solved by a minimum cut in a network with
infinite-capacity arcs on the precedence relation, a reduction due to
\citet{Picard1976} and, in the selection form, to \citet{Rhys1970} and
\citet{Balinski1970}; \citet{JohnsonNiemi1983} and \citet{Gale1957} belong to
the same early line. The polyhedral fact that makes the reduction work --- that
the constraint matrix of the closure inequalities is an interval-free network
matrix, so the natural LP relaxation is integral --- is what allows this project
to work with the LP relaxation and still speak about closures.

\citet{PicardQueyranne1980} characterise the structure of \emph{all} minimum
cuts of a network, and \citet{PicardQueyranne1982} survey the applications of
minimum cuts to combinatorial problems, including closure. Both are used here
for a specific technical point rather than for background: the distinction
between the minimal and the maximal source side of a minimum cut, which is what
the canonical convention on tied ratios needs, is theirs.

On the algorithmic side, \citet{Hochbaum2001} revisits minimum cut and maximum
flow specifically in closure graphs, \citet{HochbaumChen2000} compare
implementations on the open-pit instance of the problem, and
\citet{Hochbaum2008} gives the pseudoflow algorithm, which is the method a
serious flow-based competitor would use today. \citet{UnderwoodTolwinski1998}
is worth singling out: they solve the ultimate pit problem with a network flow
algorithm derived from the \emph{dual}, and relate its steps to those of
Lerchs--Grossmann, so the idea of attacking maximum closure through dual
information is itself old. \citet{AmankwahEtAl2014} extend Picard's reduction
to the multi-period case on a time-expanded graph, and \citet{MeagherEtAl2014}
review the design of pushbacks, which is where the sequence of nested closures
produced by the canonical decomposition is actually consumed. \citet{AhujaEtAl1993} and
\citet{Schrijver2003} are the standard references for the surrounding network
flow and combinatorial optimisation material.

A generalisation worth recording, because it places the closure problem inside a
larger family rather than beside it: \citet{HochbaumQueyranne2003} minimise a
convex cost over a closed set, which specialises to maximum closure on one side
and to isotonic regression on the other. The same family is visited from the
other end by the maximum density subgraph, which is the ratio objective of
\cref{sec:ratio} without any precedence structure at all --- solved by flow in
\citet{Goldberg1984}, given its linear program by \citet{Charikar2000}, and
treated as one parametric cut by \citet{Hochbaum2024}.

Not everything in this line is a flow algorithm, and the exception matters
enough to have its own place: the algorithm of \citet{LerchsGrossmann1965}
maintains a \emph{forest} and pivots on it. It is discussed in \cref{sec:lg},
next to the recent forest-structured work, because that is where it does its
damage.

\paragraph{What this costs the claim.} Everything about the closure model, its
LP relaxation, its integrality and its solution by minimum cut is classical.
Nothing in \FS{} may be presented as new at this level, and the project's own
documents say so: the mathematical contract states the model and attributes it,
and the tutorial recalls the reduction only to fix notation.

\section{The ratio objective}\label{sec:ratio}

Two classical devices remove a ratio from an objective. \citet{CharnesCooper1962}
rescale a linear fractional program into a linear one by normalising the
denominator, which is exactly the transformation that turns
\(\max p(S)/w(S)\) into the program with \(\sum_i w_i y_i = 1\) that \FS{}
solves. \citet{Dinkelbach1967} instead treats the ratio parametrically, solving
\(\max\{p(S) - \lambda w(S)\}\) and updating \(\lambda\) by Newton's method;
applied to closures, each iterate is a maximum-closure problem and therefore a
minimum cut. This project implements the second as one of its reference
backends.

Maximum-ratio closure --- maximum-density or maximum-mean closed set, depending
on the community --- is therefore classical in both readings, and the
equivalence between the fractional point of view and the level-set argument
(\(y = \int \chi^{C_t}\,dt\), each \(C_t\) a closure) is the standard device
relating fractional closure points to closures. The tutorial of this project
proves it in that form and attributes it there.

\subsection{The LP class, stated precisely}\label{sec:class}

Since the positioning turns on which linear program is being specialised to, it
is worth naming the class rather than the application. The program is
\begin{equation}\label{eq:class}
  \max\ \sum_{i \in R} p_i y_i
  \quad\text{s.t.}\quad
  \sum_{i \in R} w_i y_i = 1,
  \qquad y_i - y_j \le 0 \ \ \forall (i,j) \in A[R],
  \qquad \vect{y} \ge 0,
\end{equation}
with $\vect{w} > 0$, $\vect{p}$ of arbitrary sign and no upper bounds. Three equivalent
readings of the same class:

\begin{enumerate}[label=(\alph*),leftmargin=*]
\item \textbf{One budget row over a homogeneous difference system.} The
      constraints $y_i - y_j \le 0$ form a system of homogeneous difference
      inequalities on a DAG; a single dense equality row with positive
      coefficients is added.
\item \textbf{The base of an order cone.} The set $\{\vect{y} \ge 0 : y_i \le y_j\}$ is
      the cone of nonnegative order-preserving vectors on the poset, whose
      extreme rays are the characteristic vectors of the nonempty closed sets;
      the equality slices it into the polytope whose vertices are
      $\chi^C / w(C)$. Optimising over it therefore returns
      $\max_C p(C)/w(C)$, which is why the transformation of
      \citet{CharnesCooper1962} is exact here. The polytope is the cone analogue
      of Stanley's order polytope, whose vertices are the characteristic vectors
      of the filters of the poset \citep{Stanley1986}.
\item \textbf{Transposed incidence plus one dense row.} Each difference
      constraint has a $+1$ and a $-1$ in the columns of two vertex variables, so
      the matrix is the \emph{transpose} of a digraph incidence matrix, with one
      dense row appended and a right-hand side that is zero on the network part.
\end{enumerate}

Reading (c) matters for \cref{sec:forests} and was stated loosely in an earlier
draft. In the side-constrained network flow literature the variables are arcs
and the rows are vertices; here the variables are vertices and the rows are arcs. The
two are linear programming duals of one another --- the dual of
\eqref{eq:class} is $\min \beta$ subject to $w_i\beta + \divg_i(\vecg{\alpha}) \ge p_i$,
$\vecg{\alpha} \ge 0$, a flow-like system with one scalar --- so the algebra transports
between them, but pricing on one side corresponds to the ratio test on the other
and the roles are not interchangeable without care. This is the same duality
\citet{UnderwoodTolwinski1998} exploit when they attack the ultimate pit problem
through the dual. It also places \eqref{eq:class} in the same structural family
as isotonic regression and total variation on a graph, which are difference
systems on vertex variables too --- and that is why \citet{Yang2026} is closer
prior art than the network-simplex-with-a-side-row line, not further.

\paragraph{The one non-generic modelling choice.} Replacing the equality
\(\sum_i w_i y_i = 1\) by an inequality adds the origin to the feasible set and
returns value zero whenever every nonempty closure has negative profit. With
nonnegative profits this is immaterial. Here it is not: profits carry either
sign, a profitable vertex may require loss-making prerequisites, and the outer
decomposition calls the oracle on residual graphs whose profits deteriorate, so
the last calls typically have a negative optimal ratio and the decomposition
still needs the \emph{best nonempty} closure at each of them. The equality is
deliberate and is argued as such --- but it is a modelling choice inside a
classical transformation, not a contribution.

\paragraph{Recent work on the ratio side.} \citet{Hochbaum2024} gives an
incremental parametric cut method for ratio problems of this family; it is the
most recent point of comparison found for the oracle itself and is in
\file{literature/papers/}. It is a competitor to the Dinkelbach backend, not to
\FSC{}'s basis machinery, and it should be cited where the oracle is discussed.

\section{The residual decomposition and its parametric reading}\label{sec:parametric}

Repeatedly extracting a maximum-ratio set and continuing on the residual is not
a device of this project. It is the decomposition of \citet{Sidney1975} for
precedence-constrained single-machine scheduling, developed in the sequencing
literature by \citet{Lawler1978} and surveyed by \citet{CorreaSchulz2005};
\citet{MargotEtAl2003} give the version closest to the ratio reading used here:
an in-depth study of an LP relaxation of a precedence-constrained problem whose
algorithmic core is a parametric extension of Sidney's decomposition, obtained
by essentially solving a parametric minimum-cut problem.

The vocabulary of this document follows \citet{DoseEtAl2026}, the companion work
of this project, which studies the same canonical object on a precedence graph
that is a directed forest. What is called a \emph{canonical block} here is a
\emph{closure layer} there; the canonical decomposition is their \emph{optimal
sequence of closure layers}; the ratio $\lambda_r$ of a block is a
\emph{breakpoint} of $u(\lambda)$; and the per-vertex bound this project
computes is their \emph{threshold} of a vertex --- the same quantity, not an
analogy. Their scalar $\lambda$, parametrising the weights as $p_i - \lambda
w_i$, is the $\rho$ of the ratio oracle here, which is also the reduced cost
$b_i = p_i - \rho w_i$ of \cref{sec:lg}. The two works do not overlap on
content: they give aggregation algorithms specialised to forests, reaching
$\mathcal{O}(n\log n)$, whereas the precedence graph here is an arbitrary DAG
and the method is the simplex. The sequence produced is the same object, so the
two are directly comparable on instances that happen to be forests --- and on
those instances their bound is the one to beat.
The same object appears from the flow side as the sequence of nested minimum
cuts of a parametric network: \citet{EisnerSeverance1976} give the
divide-and-conquer scheme that computes all breakpoints with a number of
maximum-flow computations proportional to their number, and
\citet{GalloEtAl1989} give the parametric maximum-flow algorithm that computes
the whole path in the time of a single flow computation for monotone
parametrisations. \citet{Topkis1978} supplies the lattice/submodularity view
that makes the nesting structural rather than incidental.
\citet{PicardSmith2004} apply exactly this to open-pit contours, and
\citet{Eppstein2018} gives the modern complexity treatment of the parametric
closure problem.

This project implements the Eisner--Severance scheme as its \texttt{parametric}
backend, and its experimental report is explicit that the implementation is a
divide-and-conquer decomposition with independent maximum flows, without the
amortised flow reuse of \citet{GalloEtAl1989}.

\paragraph{What this costs the claim.} The existence of the canonical
decomposition, its nesting, the strictly decreasing ratios, the interpolation
between two consecutive prefixes that yields the LP value at a given capacity,
and the solution of the whole path by parametric flow, are all established. The
project's mathematical contract already presents them as such. A claim of
novelty at this level would be simply wrong, and the earlier drafts were right
to drop it.

\citet{NemeschEtAl2025} survey exact and approximation algorithms for
linear-parametric optimisation and place the closure case inside the general
picture; it is the reference to consult before asserting that any parametric
technique used here is unusual.

\paragraph{What survives here.} Only one thing, and it is small: the convention
that resolves ties. When several closures attain the maximum ratio, the
canonical block is their union, which on the flow side is the \emph{maximal}
source side of a minimum cut rather than the minimal one
\citep{PicardQueyranne1980,PicardQueyranne1982}. This project implements both
conventions and verifies that they agree; that is careful engineering on a known
structure, not a new result.

\section{The precedence-constrained knapsack and open-pit mining}\label{sec:pckp}

The problem the outer procedure solves is the LP relaxation of the
precedence-constrained knapsack, which has two largely separate literatures.

\paragraph{Polyhedral and exact work.} The PCKP polytope has been studied
through induced cover and clique inequalities and their lifting
\citep{Boyd1993,ParkPark1997,VanDeLeenselEtAl1999}, with separation of maximally
violated inequalities \citep{ChoShaw1997,ShawCho1998} and exact algorithms
\citep{SamphaiboonYamada2000}; \citet{KellererEtAl2004} is the standard
monograph, and \citet{KolliopoulosSteiner2007} treat the approximation side.
Tree-structured special cases have their own line
\citep{Johnson1968,Horn1972,AdolphsonHu1973}. This literature is about the
integer problem; \FS{} solves the LP relaxation, so it borders on this work
rather than competing with it. The distinction matters and the project's
documents keep it: the solver treats the LP, the integer problem is a different
problem.

\paragraph{Mining.} The ultimate-pit problem is maximum closure
\citep{LerchsGrossmann1965,Matheron1975a,Vallet1976}, and the capacitated and
multi-period versions are where the PCKP LP relaxation is actually used at
scale. \citet{BienstockZuckerberg2010} give the decomposition method for LP
relaxations of large precedence-constrained problems that is the reference point
for scale, studied further by \citet{MunozEtAl2018};
\citet{ChicoisneEtAl2012} give an algorithm for the open-pit production
scheduling problem; \citet{BolandEtAl2012} and \citet{EspinozaEtAl2015} are the
surrounding work, the latter being the MineLib benchmark this project derives
instances from. \citet{ByunDimitrakopoulos2013} solve the single-capacity
closure LP by parametric flow, which is the closest published match to the task
\FSC{} performs.

\paragraph{What this costs the claim.} Two things. First, the single-capacity
closure LP by parametric flow is already published, so \FS{} enters a solved
problem and must justify itself on how it solves it, not on what it solves.
Second, at the scale where this LP matters in practice, the incumbent methods
are flow-based and mature. The project's own experiments are consistent with
that: on the instances derived from MineLib the forest configurations solve
nothing within the time limit while the parametric backend solves all of them.

\section{Algorithms whose state is a forest}\label{sec:forests}

This is the section that decides the claim. Four bodies of work sit between
\FSC{} and any statement of the form ``an algorithm whose state is a forest on
the precedence graph''. The first of them is sixty years old and comes from
mining, which is where this project's own instances come from; it was missing
from the earlier drafts of this review, and it is the most uncomfortable of the
four.

\subsection{Normalized trees: Lerchs--Grossmann, 1965}\label{sec:lg}

\citet{LerchsGrossmann1965} solve maximum closure with a graph algorithm, not
with a flow algorithm, and the object it maintains is a forest. In the exposition
of \citet{Hochbaum2001}, from which the following is taken verbatim in substance:
the algorithm maintains a spanning tree rooted at an added vertex $r$ in the
extended graph, called a \emph{normalized tree}; each branch $T_q$ hanging from
the root carries its \emph{mass} $M_q = b(T_q)$, the sum of the vertex weights it
contains; a downward arc is \emph{strong} if the mass it supports is positive;
the tree is normalized when its only strong arcs are the ones adjacent to $r$;
and, dropping the root arcs, Hochbaum writes that ``we also refer to the forest
$(V, A_T)$ in $G$ as a normalized tree''. The invariant is that the union of the
strong branches is a maximum closed set of the forest, and the weight of the
strong vertices bounds the maximum closure value from above.

An iteration is a pivot in everything but the name. Quoting \citet{Hochbaum2001}
again: a \emph{merger} consists of ``adding an arc $(s,w)$ from a strong branch
$S$ to a weak branch $W$ and removing the arc from the root of the branch $S$ to
the root $r$'', after which the branch hangs from the strong vertex $s$ and the
supported masses are updated. One arc enters, one arc leaves, one component
loses its root, and the aggregates on the affected path are recomputed.
\citet{ZhaoKim1992} is a variant that differs in how the blocks are regrouped
after a violation, which is the split side of the same operation, and
\citet{GianniniEtAl1991} is the flow-based implementation that competed with it.

Two further facts from \citet{Hochbaum2001} matter here. First, the normalized
tree carries enough information to construct a feasible flow in linear time, so
the primal--dual correspondence between the forest state and a flow certificate
is known. Second, he gives a parametric version of the Lerchs--Grossmann
algorithm with the same complexity as a single run --- which makes it a
\emph{competitor} on the task \FS{} performs, not merely an antecedent.
\citet{PicardSmith2004} close the loop on the objective as well: they compute
intermediate open-pit contours maximising a benefit--cost ratio by parametric
maximum flow.

\paragraph{What this costs the claim.} A great deal, and more than the recent
work does. The correspondence with \FS{} is not thematic but term by term: a
branch aggregate $b(T_q)$ against the reduced cost $p(T) - \rho\,w(T)$; a strong
arc against a positive reduced cost; the union of the strong branches against
the support of the current basic solution; a merger against an entering arc with
its leaving root arc. Setting $b_i = p_i - \rho w_i$ makes the Lerchs--Grossmann
strong/weak test literally the sign of the \FS{} reduced cost. So the forest
state, the branch aggregates, the sign test on an aggregate and the merge
operation are not new relative to the 2025--2026 work: they are not new relative
to 1965. What Lerchs--Grossmann do \emph{not} do is run a simplex --- there is no
basis in the LP sense, no ratio test, no dual multiplier as an LP object --- and
they solve the closure problem at fixed weights rather than the ratio LP.
That distinction is now the whole of the surviving claim, and
\cref{sec:open}~Q5 says what has to be proved to keep it.

\subsection{Network simplex with a side constraint}

That a basis of a network LP is a spanning tree, that adding a side row makes it
a spanning structure plus one extra element, and that pivots are described as
operations on that structure, is textbook material
\citep{AhujaEtAl1993,Schrijver2003}. The paradigm ``simplex on an incidence
matrix with one side row'' is therefore not available as a contribution, and any
formulation that suggests otherwise must be rewritten.

The specific case of one side row has its own line. \citet{ChenSaigal1977} give
a revised simplex variant for capacitated network flow with additional linear
constraints in which the working basis has the size of the side system rather
than of the network. \citet{Caliskan2011} specialises the network simplex to
the constrained maximum flow problem. \citet{HolzhauserEtAl2017} give a fully
combinatorial network simplex for the budget-constrained minimum cost flow
problem, with two kinds of vertex potentials, three kinds of reduced costs,
optimality criteria proved directly on the network and an explicit anti-cycling
treatment. These are the right block to read, but they are one
transposition away from \eqref{eq:class}: their variables are arcs and their
rows are vertices, whereas here the variables are vertices and the rows are arcs
(\cref{sec:class}(c)). The two formulations are linear programming duals, so the
machinery transports and the differences have to be argued rather than assumed
(\cref{sec:open}, Q3) --- but the resemblance is a duality, not an identity, and
an earlier draft overstated it.

\subsection{Yang: a specialised simplex on rooted spanning forests}

\citet{Yang2026} gives, for budget-constrained total-variation-regularised
linear programs on graphs, a characterisation of basic solutions in terms of
rooted spanning forests with orientation, an interpretation of the simplex
method as operations on those forests, and an accelerated implementation
compared against a commercial solver. The abstract states the first two points
directly. Total variation on a graph is a difference system on vertex variables
with a budget row, so his program and \eqref{eq:class} are in the same
structural family, and not merely in the same spirit: this is the closest match
in formulation type among all the work reviewed here.

The overlap with \FS{} is substantial and was underestimated in an earlier
draft. Yang's setting has a budget row and a graph-structured constraint matrix;
when the budget slack is nonbasic he obtains a component without a root, which
plays the role of the rootless positive component here; he writes the basic
variables explicitly in terms of the nonbasic ones, derives reduced costs for
root and edge variables, states optimality in terms of tree and subtree shifts,
reads pivots as merges, splits and shifts, uses blocking edges, and updates
subtree aggregates locally.

Consequently the following claim, made in an earlier draft, is withdrawn: that
no work carries revised-simplex algebra to this level of explicitness on a
forest structure. Yang does exactly that, on a different linear program.

\subsection{Wuille and the spanning-forest linearization line}

\citet{Wuille2025} sets up, for Bitcoin transaction cluster linearization, the
program \(\max \sum_i f_i t_i\) subject to \(\sum_i s_i t_i = 1\), the dependency
inequalities and \(t \ge 0\): the same normalised program \FS{} solves, up to
orientation convention and application. The post states that the method
maintains \(n-1\) free variables, that the active dependencies form a spanning
forest, that within every chunk there is at most one free vertex variable, that
all chunks except one contain such a variable and are therefore excluded, and it
describes the pivot operations as merging and splitting chunks. These were
verified against the post itself.

\citet{MollakhaniEtAl2026} formalise that reading: the paper presents the
equivalent linear programming formulation and develops the Spanning Forest
Linearization algorithm, motivated, in the authors' words, by structural
properties of basic feasible solutions of the simplex method represented through
spanning forests. It then leaves the exact simplex: it stops maintaining the
identity of the root and a single distinguished included component, compares
chunks locally, and optimises the linearization directly.

This is the closest antecedent in absolute terms, and it removes several items
from the candidate contribution list: the normalised program itself, the count
of \(n-1\) nonbasic variables, the forest of tight slacks, the ``at most one
free vertex variable per component'', the single positive component, the
connectivity of the support, and the four qualitative pivot effects. None of
these may be presented as new.

\subsection{What the four leave open}

Lerchs--Grossmann maintain a forest on the closure graph and test the sign of
branch aggregates, but do not run a simplex and do not solve the ratio LP.
\citet{Yang2026} develops the exact simplex, but for a different program. Wuille
and Mollakhani et al.\ work on the same program, but leave the exact simplex in
favour of a direct linearization algorithm. \citet{HolzhauserEtAl2017} develops a
combinatorial simplex for a network with one side row, but on minimum cost flow.
The intersection of all four --- the ordinary primal simplex, developed exactly,
on the maximum-ratio closure program --- is where whatever remains of the
contribution has to live, and \cref{sec:positioning} states how little that is
and what would have to be shown to make it stand.

\section{Positioning: what may be claimed}\label{sec:positioning}

\subsection{Retracted claims}\label{sec:retracted}

The following statements appeared in earlier drafts of this project and are
withdrawn. They are recorded here rather than silently deleted, because a
retraction that is not written down tends to reappear in the next draft.

\begin{enumerate}[label=R\arabic*.,leftmargin=*]
\item \emph{``First complete closed-form realization of the two revised-simplex
      basis systems for a forest-structured graph LP.''} Withdrawn: too close to
      \citet{Yang2026}, who develops exactly that on a budget-constrained
      total-variation program.
\item \emph{``First closed-form forest simplex.''} Withdrawn: too broad, and
      contradicted both by \citet{Yang2026} and by the classical
      network-simplex-with-side-constraint literature
      \citep{AhujaEtAl1993,Schrijver2003}.
\item \emph{Novelty of the forest structure of basic feasible solutions.}
      Withdrawn: present in \citet{Wuille2025} and formalised in
      \citet{MollakhaniEtAl2026} for the same normalised program.
\item \emph{Novelty of the merge/split reading of pivots.} Withdrawn: same
      sources, and \citet{Yang2026} independently.
\item \emph{Novelty of the forest state and of the aggregate sign test, even
      relative to the 2025--2026 work.} Withdrawn, and further back than the
      previous two items: \citet{LerchsGrossmann1965} maintain a forest of
      branches on the closure graph, carry the mass of each branch, classify an
      arc by the sign of the mass it supports, and pivot by adding one arc and
      removing the branch's root arc (\cref{sec:lg}). With $b_i = p_i - \rho
      w_i$ their strong/weak test is the sign of the \FS{} reduced cost. Nothing
      about maintaining a forest and testing aggregates on it can be claimed.
\item \emph{Novelty of the maximum-ratio closure formulation, of its solution by
      max-flow, and of the canonical decomposition.} Withdrawn: classical
      \citep{Picard1976,PicardQueyranne1982,Sidney1975,Lawler1978}.
\item \emph{Novelty of the ``simplex on an incidence matrix plus one side row''
      paradigm.} Withdrawn, and datable. \citet{ChenSaigal1977} prove that a
      basis of a capacitated flow problem with $k$ additional linear rows is a
      spanning tree plus $k$ columns, and drive pricing, pivoting and inversion
      from a working basis of size $k$; \citet{Caliskan2011} specialises the
      case $k=1$ to maximum flow with a budget row, where the basis is a
      spanning tree plus one arc closing a fundamental cycle of nonzero net
      cost; \citet{HolzhauserEtAl2017} do the same for minimum cost flow. The
      paradigm is due to 1977, not to the recent literature.
\item \emph{Novelty of solving the ratio version by walking a forest state.}
      Withdrawn: \citet{Hochbaum2001} gives a parametric Lerchs--Grossmann with
      the complexity of a single run, and \citet{PicardSmith2004} compute
      open-pit contours maximising a benefit--cost ratio by parametric flow.
\item \emph{Novelty of the equality normalisation \(\sum_i w_i y_i = 1\).}
      Withdrawn: it is the homogeneous case of the reduction of
      \citet{CharnesCooper1962}, whose Theorem~2 already writes the
      normalisation as an equality with right-hand side $1$
      (\cref{sec:open}, Q5).
\item \emph{Any claim of superior practical performance.} Withdrawn on the
      project's own evidence: on the frozen test bed the parametric and
      Dinkelbach max-flow backends dominate \FS{} on both solved count and time
      (\cref{sec:open}).
\end{enumerate}

\subsection{The cell nobody occupies}\label{sec:gap}

The claim reduces to one empty cell in a table with two axes: what the method
maintains, and which program it solves.

\begin{center}
\begin{tabular}{@{}lll@{}}
\toprule
& maximum closure, weights fixed & maximum-ratio closure \\
\midrule
flow / cut state
  & \citet{Picard1976}, \citet{Hochbaum2008}
  & \citet{GalloEtAl1989}, \citet{Hochbaum2024} \\
forest state, pivots, no LP basis
  & \citet{LerchsGrossmann1965}
  & \citet{Hochbaum2001} (parametric LG), \citet{MollakhaniEtAl2026} \\
exact simplex, LP basis
  & ---
  & \textbf{this work} \\
\addlinespace
\multicolumn{3}{@{}l@{}}{\emph{exact simplex on a sibling program:}
  \citet{Yang2026} (budget-constrained total variation),
  \citet{HolzhauserEtAl2017} (budget-constrained min cost flow),} \\
\multicolumn{3}{@{}l@{}}{\quad
  \citet{ChenSaigal1977} ($k$ side rows on a capacitated flow),
  \citet{Caliskan2011} (budget-constrained max flow)} \\
\bottomrule
\end{tabular}
\end{center}

Read across the bottom row: the exact simplex has been developed on programs
in the same structural family as \eqref{eq:class} --- a difference system or a
network with one budget row --- but not on \eqref{eq:class} itself. Read up the
right column: the ratio program has been solved by flow states and by forest
states that pivot, but not by a simplex maintaining an LP basis with reduced
costs, a ratio test and dual multipliers. \citet{MollakhaniEtAl2026} say so
themselves: their algorithm is, in the abstract's own words, motivated by
structural properties of basic feasible solutions of the simplex method, and
what it then maintains and refines is a partition of the vertices into
dependency-respecting subsets ordered by decreasing aggregate ratio --- chunks
merged and split directly on the graph, not a basis.

\paragraph{The most serious candidate for the empty cell, and why it does not
fill it.} \citet{MargotEtAl2003} is the one paper found in this review that
promised to occupy the bottom-right cell: an in-depth theoretical, algorithmic
and computational study of an LP relaxation of a precedence-constrained problem,
whose algorithmic core is a parametric extension of Sidney's decomposition and
whose finest decomposition is obtained, in their words, by essentially solving a
parametric minimum-cut problem. Read in full, it belongs in the first row and
not in the third: the state it maintains is a sequence of nested minimum cuts,
and the only simplex in the paper is the one inside the general-purpose LP
solver used to produce the computational results. It carries no basis
characterisation, no reduced costs and no ratio test for the relaxation. It
strengthens R6 and R8 --- the decomposition and its parametric reading are
classical --- and leaves the cell empty.

\paragraph{A warning that belongs next to the table.} An empty cell is not by
itself a contribution. It can be empty because the route is worse, and the
evidence available here points that way: on this project's own test bed the
flow-based entries of the first row and the parametric entry of the second
dominate the bottom row on both solved count and time (\cref{sec:open}). The
honest reading is that this work fills the cell and reports what is inside it,
including the part that is unfavourable.

\subsection{What is left}

Four candidates survive the review. None of them is a paradigm; all of them are
specialisations, and each carries a condition that has not yet been discharged.

\begin{enumerate}[label=C\arabic*.,leftmargin=*]
\item \textbf{Specialisation of the ordinary simplex, primal and dual, to the
      normalised maximum-ratio closure LP, and to the canonical decomposition
      it drives.} That a simplex method is exact is not a contribution and is
      not offered as one; what is at stake is what the method \emph{becomes}
      on this program. \citet{Wuille2025} and
      \citet{MollakhaniEtAl2026} expose the forest structure of the basic
      feasible solutions of this program and then leave the simplex, optimising
      the linearization directly; \citet{Yang2026} develops the exact simplex,
      but on a different program; \citet{LerchsGrossmann1965} pivot on a forest
      of the closure graph, but outside the simplex and at fixed weights.
      Carrying the ordinary simplex through on this LP --- basis state, prices,
      pivot directions, ratio test, multipliers, certificates, in a primal and
      in a dual variant --- and then driving the whole canonical sequence of
      blocks with it, not merely the first, is the gap between all three.
      \emph{Condition:} it must be presented as a specialisation, not as a new
      method; the debt to all three lines must be stated in the abstract, not
      only in related work; and the difference from a normalized tree must be
      made precise rather than asserted, because to a mining audience the two
      objects look identical.

\item \textbf{The closure-specific closed forms are markedly simpler than the
      general forest-simplex algebra.} For the normalised program the reduced
      costs collapse to
      \[
        \bar c_{y_T} = p(T) - \rho\,w(T),
        \qquad
        \bar c_{s_e} = \sigma_e\bigl[p(D_e) - \rho\,w(D_e)\bigr],
      \]
      with no edge objective coefficients, no box-bound contributions, no
      orientation choice between positive and negative edge variables and no
      generic working-basis algebra. \emph{Condition:} the comparison against
      \citet{Yang2026} must be made explicitly and term by term; ``simpler''
      is not a claim until the two sets of formulas are put side by side.

\item \textbf{Full pricing in $O(n)$ arithmetic, independent of the number of
      precedence arcs.} There are exactly $n-1$ nonbasic candidates; all forest
      aggregates are obtained in one postorder; each candidate costs $O(1)$; no
      non-forest arc is ever priced. \emph{Condition:} this is the cleanest of
      the four, and also the one most exposed. If \citet{Yang2026} admits an
      equivalent $O(n)$ pricing in his model, the claim is not that $O(n)$
      pricing is possible but that it is available here without the structural
      overhead his model carries. That check has not been done.

\item \textbf{The separation between pricing and feasibility in the exact ratio
      test.} Pricing ranges over the $n-1$ nonbasic forest and root variables;
      feasibility ranges over all potentially decreasing basic slacks. The
      distinction is clean and it is what makes the exact simplex correct rather
      than merely plausible. \emph{Condition:} on its own this is a component of
      C1, not an independent contribution, and should be presented as such.
\end{enumerate}

\subsection{The formulation that survives}

\begin{quote}
To the best of our knowledge we give the first explicit simplex realisation,
primal and dual, specialised to the normalised maximum-ratio closure LP, and we
use it to produce the entire canonical sequence of blocks rather than the
maximum-ratio block alone. Building on the forest structure of its basic
feasible solutions --- which is closely related to recent spanning-forest
linearization work --- we derive problem-specific closed forms for all nonbasic
prices and pivot directions. In particular, complete pricing requires $O(n)$
arithmetic independently of the number of precedence arcs, while a full ratio
test accounts for all potentially blocking basic slacks. The canonical
decomposition itself is classical and is not claimed.
\end{quote}

\noindent followed immediately, and in the same place, by:

\begin{quote}
We do not claim novelty for the forest structure itself, for the merge/split
basis effects, for the maximum-ratio closure formulation, for the canonical
decomposition, or for the general paradigm of simplex methods on incidence
matrices with side constraints; these have direct precedents in
\citet{Wuille2025} and \citet{MollakhaniEtAl2026}, in \citet{Yang2026}, and in
the classical network-simplex literature.
\end{quote}

\subsection{The name}\label{sec:name}

The name follows the positioning, and the earlier one contradicted it. ``Forest
Simplex'' named the forest --- the one element \cref{sec:retracted} withdraws
--- and it collided three ways: with the \emph{spanning forest} of
\citet{MollakhaniEtAl2026} and \citet{Wuille2025}, with the \emph{rooted
spanning forests} of \citet{Yang2026}, and, one letter and no syllables away,
with the Forrest--Tomlin update of the basis factors, which is what an audience
of linear programmers hears first.

The method is therefore called the \FS{}, lower case, on the model of
\emph{network simplex}: a name taken from the problem the method is specialised
to, not from the structure it maintains. The problem is
\(\max_C p(C)/w(C)\) over the closures of a digraph, and the linear program is
\eqref{eq:class}; the capacity row of the application is absent from both,
which is why it is absent from the name as well. The precedence-constrained knapsack is where
the program comes from and where the method is measured, and it is reported as
such among the applications; it is kept out of the title as well as out of the
name, because the method delivers the whole canonical sequence and not only the
first block, and a knapsack in the title would promise that the capacity row is
handled by the simplex, which it is not.
The forest keeps its own name, as the structure of the basis, with the
attribution \cref{sec:lg} gives it.

\subsection{The question the work answers}

The question is not whether maximum-ratio closure can be solved with forests ---
it can, and that is known. It is:

\begin{quote}
what does the exact ordinary primal simplex look like on the normalised
maximum-ratio closure LP, once the known forest structure of its basic feasible
solutions is exploited completely?
\end{quote}

The answer is the algorithm: basis state as forest plus root, exact aggregate
prices, $O(n)$ pricing, exact forest motion, full ratio test, dual multipliers
and certificates, and an implementation measured against the alternatives. That
last item is where the work currently stands weakest, and \cref{sec:open} says
why.

\section{What has to be settled before claiming}\label{sec:open}

\subsection{The checks this review could not close}

\begin{enumerate}[label=Q\arabic*.,leftmargin=*]
\item \textbf{Term-by-term comparison with \citet{Yang2026}.} Contribution C2
      asserts that the closure-specific closed forms are markedly simpler than
      the general forest-simplex algebra, and C3 asserts that $O(n)$ pricing is
      available here without the structural overhead of Yang's model. Neither
      has been verified against the paper's actual formulas. This is the single
      most load-bearing open item: if Yang's development specialises to
      something equivalent, C2 and C3 both collapse into ``an instance of a
      known scheme''.
\item \textbf{Exact delimitation against \citet{MollakhaniEtAl2026}.} Their
      abstract settles the direction: the algorithm is motivated by the
      structure of basic feasible solutions of the simplex method, and it then
      merges and splits chunks directly on the dependency graph, maintaining a
      partition ordered by decreasing aggregate ratio rather than an LP basis.
      That is a departure from the simplex, not a specialisation of it. What
      remains unchecked is the full text: whether the paper anywhere derives
      reduced costs, a ratio test or dual multipliers for \eqref{eq:class}, and
      how far its optimality argument coincides with the pricing argument used
      here.
\item \textbf{The network simplex with one side row.} \emph{Closed on the
      structural question; open only on the term-by-term comparison.}
      \citet{ChenSaigal1977} and \citet{Caliskan2011} have now been obtained
      and read. They push this block back to 1977 and settle what it does and
      does not cover.

      \citet{ChenSaigal1977} treat a capacitated flow problem with $k$
      arbitrary additional linear rows. Their Proposition 1.1.5 is the
      structural statement: in any partition of a basis, the square block
      corresponds naturally to a \emph{spanning tree} of the network, so a basis
      is a spanning tree plus $k$ further columns, and the subnetwork carried by
      the basis has rank $n-1$ and nullity $k$. Pivoting, pricing and inversion
      are then driven by a \emph{working basis} of size $k$ rather than $n+k-1$.
      \citet{Caliskan2011} is the case $k=1$ made explicit on maximum flow with
      a budget row: a basis is the two subtrees rooted at $s$ and $t$, the arc
      $(t,s)$, and \emph{one} further arc closing a fundamental cycle, with his
      Lemma~1 recording that the net cost around that cycle cannot vanish, which
      is the rank condition the side row imposes.

      Two consequences. First, R7 in \cref{sec:retracted} is confirmed and can
      now be cited precisely instead of being called textbook: tree-plus-$k$
      bases, a working basis of the size of the side rows, and a specialised
      pricing built on them are due to 1977, not to 2017. Second, the phrase
      ``to the best of our knowledge'' in \cref{sec:positioning} survives, for a
      reason \cref{sec:ratio} already states: this whole block works on the
      node--arc incidence matrix with variables on \emph{arcs}, where a basis is
      a spanning tree, while \eqref{eq:class} carries variables on
      \emph{vertices} and difference constraints on arcs --- the
      \emph{transposed} incidence matrix --- where a basis is a forest with
      anchored roots and not a spanning tree of the same graph. The paradigm is
      theirs; the specialisation to the transposed structure is not in any of
      the three papers.

      What remains open is narrower and is the same check as Q1:
      \citet{HolzhauserEtAl2017}, the closest of the block, has been obtained
      but not read term by term, so the claim that the closure-specific closed
      forms are simpler than the generic working-basis algebra is still asserted
      rather than demonstrated for that paper.
\item \textbf{Formal separation from the normalized tree.} This is now as
      load-bearing as Q1. A normalized tree of \citet{LerchsGrossmann1965} and
      an \FS{} basis are, on inspection, the same object read twice: forest
      plus root, aggregates per branch, a sign test per branch, one arc in and
      one root arc out. What has to be written down is the exact statement of
      the correspondence and of where it breaks --- presumably that a normalized
      tree is not a basis of the ratio LP, that it carries no ratio test and no
      dual multiplier, and that the parametric Lerchs--Grossmann of
      \citet{Hochbaum2001} walks breakpoints rather than performing simplex
      pivots on the normalised program. If instead the correspondence turns out
      to be exact, then \FS{} is a reinterpretation of a 1965 algorithm, which
      is a publishable observation but a completely different paper from the one
      currently drafted.
\item \textbf{Prior art on the equality normalisation.} \emph{Closed, and it
      is a precedent.} \citet{CharnesCooper1962} has now been obtained and
      read. Their Theorem~2 reduces a linear fractional program to the ordinary
      linear program
      \[
        \max\ c^{\mathsf T}y + \alpha t
        \quad\text{s.t.}\quad
        Ay - bt \le 0,\qquad
        d^{\mathsf T}y + \beta t = 1,\qquad y,t \ge 0,
      \]
      in which the normalisation is an \emph{equality} and the constant is
      exactly $1$. The program \eqref{eq:class} that \FS{} solves is the
      homogeneous case $\alpha = \beta = 0$, $d = w$, of that transformation.
      The deliberate choice of an equality is therefore not a modelling
      contribution of this project; it is the 1962 construction, and
      \cref{sec:ratio} already cited it as such.

      Their treatment of sign is worth recording because it is \emph{not} ours.
      Charnes and Cooper require $\operatorname{sgn}(\gamma) =
      \operatorname{sgn}(d^{\mathsf T}x^\ast + \beta)$ and therefore solve two
      linear programs, the second obtained by negating $(c,\alpha)$ and
      $(d,\beta)$ --- ``a change in sign in the functional and in one
      constraint''. Here the weights are positive, so the denominator is
      positive on the whole feasible cone and only their first program is ever
      needed, while a negative \emph{numerator} is representable inside it.
      That is an observation about the instance class, not a device, and it
      should not be presented as one.
\end{enumerate}

\subsection{The computational axis}

The project's own experimental evidence does not support a performance claim.
On the frozen test bed the max-flow backends --- parametric decomposition and
Dinkelbach --- solve more instances, and solve them faster, than \FS{} in every
configuration tried; the numbers are in the experimental report and are not
reproduced here, because they are being re-measured on a quiet machine and this
document must not carry a second, drifting copy of them. The qualitative
finding is stable across every campaign run so far and is unlikely to move: the
dominant cost is degeneracy, the primal spends the overwhelming majority of its
pivots on zero steps, and the dual removes the degeneracy without buying a
comparable speedup, trading null steps for small ones.

The consequence for positioning is that the work is, today, a
\emph{specialisation} paper and not a \emph{performance} paper, and the three
experiments that could change that have not been run:

\begin{enumerate}[label=E\arabic*.,leftmargin=*]
\item \textbf{Single ratio oracle.} \FS{} against a spanning-forest
      linearization implementation adapted to the first-chunk task, against a
      static min-cut/root search, and against a general LP solver.
\item \textbf{Full decomposition path.} Repeated residual oracles against
      spanning-forest linearization, against parametric max-flow
      \citep{GalloEtAl1989}, against the parametric Lerchs--Grossmann of
      \citet{Hochbaum2001}, and against a mining-style implementation
      \citep{ByunDimitrakopoulos2013} where available.
\item \textbf{Sequences of related instances.} Vertex deletions, forced
      inclusions and exclusions, residual prefix removal, small profit and
      weight perturbations, arc additions and removals --- comparing forest
      basis reuse against pseudoflow warm start, against parametric state
      reuse, and against cold restart.
\end{enumerate}

E2 was the natural place to look for an advantage, and it has now been measured
on the full canonical sequence rather than on a single capacity. The result is
negative and the reason is instructive, so both are recorded here.

The task \texttt{full-path} was run on twelve instances spanning the generated
families, a historical PCKP instance and a MineLib instance, comparing the tuned
primal, the dual, Dinkelbach and the parametric decomposition, two replicas
each, under a sixty-second limit. The parametric backend is fastest on eleven of
the twelve; the dual wins only on the arc-free graph, where there is nothing to
pivot on. The margin is not marginal: the median ratio between the tuned primal
and the parametric backend is a factor of twelve, and on a layered DAG with $n =
2000$ it is a factor of $859$ --- sixty seconds against seven hundredths. On the
largest instance tried, a sparse DAG with $n = 20{,}000$, the tuned primal, the
dual and Dinkelbach all exhaust the limit while the parametric backend finishes
the whole sequence in $3.6$ seconds.

The diagnosis matters more than the ranking. The oracle count is exactly what
the design predicts --- one ratio oracle per closure layer, the same count
Dinkelbach pays --- and the parametric backend uses about two flow computations
per layer. The difference is entirely inside the oracle: each one costs tens to
hundreds of thousands of pivots, of which between $98\%$ and $100\%$ are
degenerate. On the layered instance the primal spends $768{,}639$ pivots, essentially
all of them zero steps, to produce a sequence the parametric backend obtains
with a few hundred flow computations.

The warm start was then implemented and measured, and it recovers most of that
cost. The extracted block \emph{is} the positive component, that is one tree of
the basis forest, so the other trees survive untouched and stay anchored: the
next oracle can reuse that forest instead of rebuilding from the empty basis.
The condition that was missing from the project's earlier decision is a single
one and is checked rather than proved: the values are $1/w(T)$ on the positive
tree $T$ and zero elsewhere, so an active arc leaving $T$ would violate
feasibility and $T$ must be closed in the residual. The best closed surviving
tree is chosen, and if none is closed the oracle restarts cold. The check costs
$\mathcal{O}(m)$ per oracle.

The effect is on the pivot count, which is where the time was: $251{,}780$
pivots become $3{,}080$ on a disconnected DAG, $89{,}468$ become $6{,}108$ on a
sparse one, $48{,}974$ become $6{,}915$ on \texttt{minelib\_newman1}. The
decomposition produced is identical block by block and the objective of a
capacity query is unchanged to nine decimals. In wall time the warm primal is
$3.1$ times faster than the cold one, overtakes Dinkelbach, and wins three of
the ten instances every configuration solves --- the families whose canonical
sequence is longest --- against one of twelve before.

This changes the size of the gap, not its direction. The parametric backend
still wins the other seven instances and is $3.4$ times faster in total, so the
positioning of \cref{sec:positioning} stands: this is a specialisation paper,
not a performance paper. What has changed is that the axis is no longer
hypothetical. The degeneracy is reduced, not removed, and two measurements are
still owed: the frozen configuration was tuned on single-capacity solves and
never on the full decomposition --- a paired ablation on this task puts a sink
initial basis $1.21$ times ahead of it --- and the anti-cycling regime has never
been tuned for this task at all.

An end-to-end
branch-and-cut comparison comes after that, not before, and must not be used as
the first piece of evidence.

\subsection{Honest summary}

For the forest structure, the basic-solution characterisation, the merge/split
pivots, the sign test on branch aggregates, maximum-ratio closure, the canonical
decomposition and the incidence-matrix-plus-side-row paradigm, \FS{} is not new
--- and the forest itself is not new since 1965, not since 2025. What may be new is
the specialisation of the ordinary simplex, primal and dual, to this particular
linear program, the closed forms that specialisation produces, and the use of it
to drive the entire canonical sequence of blocks rather than the first ---
and that is a narrow claim which \cref{sec:positioning} states narrowly, and
which Q1, Q2 and Q4 have to close before it is made in public --- Q3 and Q5
having been closed by the readings recorded above, both of them against the
claim rather than for it.


\bibliographystyle{plainnat}
\bibliography{references}
\end{document}
